\chapter{Experimental Methodology}\label{ch:methodology}

\needswork{STATUS: DRAFTED from \texttt{docs/04}. The statistical rules below are binding on every
result in \Cref{ch:bytes}--\Cref{ch:lowrate}.}

\section{Binding statistical rules}

Every experiment obeys the following, and where a result predates a rule that is recorded rather
than retrofitted:

\begin{itemize}
\item At least $30$ seeded repetitions per configuration, with seeds logged in the artifact.
\item Report median and bootstrap $95\%$ confidence interval over $10^4$ resamples.
\item Microbenchmarks: $\geq 10^4$ iterations after $10^3$ warm-up, garbage collection disabled,
      runs spanning a suspend discarded.
\item Comparisons state effect size; never bare means.
\item Raw artifacts are immutable once frozen.
\item Any model-versus-measurement deviation above $10\%$ must be investigated and either explained
      in writing or the model corrected. \textbf{Silent tolerance widening is forbidden.}
\end{itemize}

\subsection{Why 30 seeds, and what it does not buy}

The 30-seed rule is not arbitrary and is not original policy --- it was imposed \emph{after} four
headline numbers were found to have been distorted by small-sample means evaluated against
thresholds. None was a modelling error; every one was sampling.

\begin{remark}[The limit of the rule]
More seeds fixes exactly one class of defect: a small-sample mean landing on the wrong side of a
threshold. It does nothing for a threshold applied to a mean rather than a distribution, for an
unverified constant on the measurement path, for a configuration change that perturbs the random
realisation, or for a claim wider than the experiment supports. \Cref{ch:reproducibility}
catalogues instances of each, several found in this project's own work \emph{after} the 30-seed
rule was in force.
\end{remark}

\subsection{Thresholds must be applied to distributions}

A criterion of the form $V \geq 0.95$ is a \emph{reliability} requirement. Applying it to a mean
across seeds answers a throughput question instead. Where this thesis reports a threshold crossing
it reports the distribution alongside --- minimum, maximum, standard deviation, and the count of
seeds that individually fail --- and where the mean and per-realisation readings disagree, both are
given.

\section{Pre-registration}

Where a run tests a prediction, the prediction is committed to version control \emph{before the
data exist}, as its own commit, so the ordering is third-party checkable in the repository history.

This practice is enforced rather than aspirational, and it has cost the project a claim: one
pre-registration assertion was \emph{withdrawn} rather than reconstructed, because writing an
expectations file after results are known manufactures evidence. It has also exposed a subtler
defect in this author's own practice, recorded in \Cref{ch:validation}: a pre-registered
\emph{directional} prediction that the experiment had no statistical power to resolve.

\needswork{TO ADD: a table of every pre-registration in the project with its commit hash, the
prediction, and the outcome. The material exists in \texttt{docs/}; the table does not.}

\section{Experiment matrix}

\begin{center}
\begin{tabular}{lll}
\toprule
Exp & Question & Primary output \\
\midrule
E1 & overhead dominance (T1) & bytes/record, $\varphi$ per encoding \\
E2 & batching cure (T2/T2a) & $\varphi$ and bytes/rec vs $b$; binding-ceiling check \\
E3 & loss frontier (T3) & measured $V$ vs $(1-p)^n$; goodput; Pareto plot \\
E4 & scheme selection (T4) & energy/record, verification wall, crossover \\
E5 & co-design end-to-end (T5) & bytes/energy/goodput vs baselines \\
E6--E7 & low-rate arm and its co-design & feasibility partition, Pareto front \\
E8 & capacity envelope & channel utilisation, $\Nmax$ \\
\bottomrule
\end{tabular}
\end{center}

Baselines throughout are A+JSON (naive), A+CBOR (compression-only) and D over-aggregated.

\section{Simulation campaign}

The simulation scenario uses IEEE 802.11a ad-hoc at $\SI{6}{\mega\bit\per\second}$, RTS/CTS off,
with $N$ saturated senders in \emph{both} unicast and broadcast modes, frame sizes taken from the
real framer output. Transport is a packet socket rather than a flow monitor, because MAC-level
goodput with no addressing artifacts is what the models predict.

\begin{remark}[The comparison rule]
Each mode is compared \emph{only} to its matching analytic variant: unicast to the acknowledged DCF
model, broadcast to the published broadcast model. Mixing them is the single most effective way to
fabricate a model gap, and this project did exactly that for several phases before catching it.
\end{remark}

\section{Hardware campaign}

Two Raspberry Pi 4 units, headless with the performance governor pinned. Timings are re-run from
the microbenchmark suite; energy uses external current sensors on a separate measurement host with
a GPIO synchronisation line, so the benchmarked device is not instrumented by itself.

The energy harness prints its \emph{predicted} value before reading the meter, deliberately, so a
measurement cannot be quietly fitted to expectation.

\needswork{TO ADD: the measurement uncertainty budget for the energy rig. The residual model gap is
stated in \Cref{ch:validation} but its error bars are not derived.}
